\BeleskeID{02-s012}
% element: 02-s012:e01
% element: 02-s012:e02
% element: 02-s012:e03
% element: 02-s012:e04
% element: 02-s012:d01

Red \(bb1\) obuhvata kombinacije 001, 011, 101 i 111; red \(b1b\) obuhvata 010, 011, 110 i 111; red \(1bb\) obuhvata 100, 101, 110 i 111. Ovi redovi se delimično preklapaju, ali u svakom zajedničkom slučaju zahtevaju isti izlaz, 1. Zajedno sa zasebnim redom 000 obuhvataju svih osam kombinacija.

Ovde je „bilo šta“ oznaka \emph{ulazne} pozicije: kada je jedan ulaz već 1, preostali ulazi ne utiču na rezultat ILI funkcije. Izlaz tog skraćenog reda je određen. To treba razlikovati od neodređenog \emph{izlaza} u zadatku minimizacije, kod koga će biti dozvoljen izbor vrednosti funkcije.
